\section{Related Work}


\emph{Systolic Arrays and Spatial Accelerators.}
Extensive work has explored the automatic generation of spatial accelerators for
target applications~\cite{polysa,autosa,dsagen,soda,interstellar,lego}.
These approaches require the application to be written in a specific form and
generate application-specific hardware accelerators.
In contrast, \MSAGA{} is compatible with existing general-purpose systolic arrays
for matrix multiplication.
Stellar~\cite{stellar} and Gemmini~\cite{ucb-gemmini} expose ISAs for
programmability, but they do not support the \FA{} algorithm.
PICACHU~\cite{picachu} proposes a CGRA-based solution to accelerate nonlinear
operations in Transformer models, including the \texttt{exp} in \FA{}.
However, it instantiates a separate array alongside the main MAC systolic array,
increasing area and incurring additional data movement overhead.


\emph{Fusing \FA{} on Systolic Arrays.}
Significant efforts have been made to fuse \FA{} on hardware accelerators. 
COSA~\cite{COSA} and COSA Plus~\cite{COSAPlus} fuse \FA{} into two 
cooperative systolic arrays, where data from the first array must pass 
through a special function unit (SFU) for softmax. Matching the throughput of SFU 
and large systolic arrays can increase hardware cost.  
\FuseMax{}~\cite{fusemax} builds on the cascade Einsum abstraction 
introduced in TeAAL~\cite{teaal} to describe the \FA{} algorithm 
and map it onto spatial accelerator models. 
\FuseMax{} uses a different dataflow than ours: by employing an output-stationary dataflow,
it overlaps four \FA{} 
iterations on the systolic array to obtain high hardware utilization.
However, the overlapping of iterations necessitates storing multiple input tiles in SRAM and
maintaining intermediate results within the array, leading to higher storage overhead 
and control complexity.
Additionally, \FuseMax{} assumes all \FA{} operations are performed in 
\texttt{fp16} format. By contrast, the original \FA{} algorithm uses 
\texttt{fp16} for matrix multiplications but \texttt{fp32} for accumulations. 
Reconciling this difference may require intermediate results to be stored in 
\texttt{fp32}, which can further increase register usage.
ExpMul~\cite{expmul} fuses \texttt{exp} and the second matrix multiplication in \FA{}.
Its approach is orthogonal to ours and could be integrated into \MSAGA{}.
